\section{Introduction}

% Present the background to your project, why the problem that you are addressing is interesting and challenging from the perspective of responsible design of interactive AI systems. Refer to the theories presented in the course that you expect to be particularly relevant in this project. Refer also to research literature with a reference list that shows the state-of-the-art of how similar problems have been addressed.

Mathematics classrooms face a persistent challenge: a typical class of 10 to 12 year-olds often contains students ranging from a primary grade level to a lower secondary level, creating a situation where high performers frequently become bored, while low performers experience anxiety and helplessness~\cite{tabassum2024managing}. Traditional teaching methods often struggle to address these individual needs simultaneously.

The recent emergence of Large Language Models (LLMs) offers new possibilities for personalized tutoring, and as such, research and implementation have shifted from \enquote{one-size-fits-all} methods towards AI-driven personalized approaches. Models such as Personalized Adaptive Gamified e-Learning (PAGE) use machine learning to process student data and dynamically adjust game mechanics, content difficulty, and feedback based on real-time progress, delivering customized learning paths aligned with individual cognitive profiles~\cite{maher2020learners}. At the modeling level, Reinforcement Learning (RL) and user modeling techniques construct representations of a learner's goals, knowledge, and misconceptions, then optimize difficulty levels and reward structures to keep learners in a productive zone between boredom and frustration~\cite{kobsa1990user}.

While these tools offer significant potential for personalization, current implementations often present new risks to the learning process~\cite{pt2025beyond}. Students frequently utilize AI tools to obtain immediate answers rather than to understand underlying concepts, leading to a superficial \enquote{illusion of knowledge} where material is not fully internalized. Furthermore, standard LLMs can suffer from \enquote{hallucinations}, producing confident but incorrect mathematical reasoning, which undermines trust in educational settings~\cite{steyvers2025what}. To address these issues, educational AI must shift from being an oracle that delivers answers to a \enquote{More Knowledgeable Peer} that facilitates Active Learning, encouraging students to internalize material through interactive engagement and constructive problem-solving.

Making this shift concrete requires grounding the system in how learning actually works. Research on cognitive agency suggests that reasoning skills develop best when learners actively engage with problems rather than receive ready-made answers~\cite{nirenburg2017cognitive, ricci2022go}, and Self-Determination Theory identifies autonomy, competence, and relatedness as the psychological conditions under which this kind of intrinsic engagement is most likely to occur~\cite{ryan2000self}. At the instructional level, the "MiniGuide" strategy provides a structured approach: rather than resolving a student's difficulty directly, the tutor poses guiding questions that keep the learner doing the cognitive work~\cite{lithner2025miniguide}.